
This section details the experimental design for the study, outlining the complete process of operationalizing the coupled ABM-RL model to answer the core research questions. This framework was the "experimental" phase of the research, defining how the simulation was run and how the results were measured. It was divided into four key stages: (1) the protocol for initializing, calibrating, and training the model; (2) the design of the specific policy scenarios to be tested; (3) the performance metrics that were used to evaluate and compare the outcomes of each scenario; and (4) the sensitivity analysis to validate the model's robustness.

\begin{figure}[htbp]
    \centering
    \includegraphics[width=0.9\linewidth]{abm-drl_framework.png}
    \caption{Process Flow of Proposed Methodology}
    \label{fig:process_flow}
\end{figure}

This diagram illustrated the integrated simulation pipeline designed to optimize waste management policies. The workflow initiated with Input \& Parameterization, grounding the model in real-world demographic data from the Philippine Statistics Authority (PSA), waste management records from the Bacolod LGU, and behavioral parameters synthesized from literature. These inputs drove the ABM \& DRL Dynamic Loop, a cyclical process where the Reinforcement Learning agent (LGU) iteratively adjusted policy levers—specifically IEC intensity, incentives, and enforcement—based on the feedback received from Household Agents. The system generated quantitative Output Results, such as segregation rates and budget utilization, which were finally translated into actionable Policy Suggestions for the Municipality of Bacolod to effectively implement R.A. 9003.

\subsection{Training and Validation}
The simulation framework was operationalized through a three-stage process: initialization, calibration, and large-scale training. 

First, the model was initialized using baseline data gathered from key-informant interviews with officials from all 7 barangays, loading in their specific population counts, local budgets, and current compliance rate estimates \citep{Villanueva2021, Yazawa2025}. This step was critical to ensure the Agent-Based Model (ABM) accurately captured the real-world heterogeneity and specific challenges facing each local unit, providing a relevant starting point for the policy simulation \citep{Camarillo2021}.

Second, the model underwent a rigorous calibration phase where the synthesized behavioral parameters—such as household sensitivity to fines and the strength of social norms—were systematically adjusted \citep{Taraghi2025, Ceschi2021}. The objective was to tune these unobservable parameters until the model's baseline behavior (when run without any new policy interventions) accurately reproduced the current-day, real-world compliance estimates provided by the LGU (e.g., the sim 10\% municipal rate) \citep{Jimenez2025}. This process of matching model outputs to observed real-world outcomes was essential for validating the ABM as a credible representation of the socio-environmental system before proceeding to policy testing \citep{Ma2023}.

Finally, once the ABM was calibrated, the Heuristic-Guided Deep Reinforcement Learning (HuDRL) agent's training began. This involved running over 12 simulated episodes, or "lifetimes" (e.g., 10,000 simulated 3-year periods), a high volume of interaction necessary to learn an optimal policy \citep{Dey2025}. By serving as the training environment, the ABM allowed the LGU agent to thoroughly explore the vast and complex 21-dimensional policy space through trial and error—a process that would have been impossible or unethical in the real world—to ultimately converge on a stable and optimal budget allocation strategy that maximized long-term compliance under fiscal constraints \citep{TianReview2024}.

\subsection{Policy Scenario Analysis}
The approach of testing distinct, constrained policy scenarios was a standard and necessary step in Deep Reinforcement Learning (DRL) studies embedded within Agent-Based Models (ABM), as it allowed researchers to isolate the causal impact of different instruments on system outcomes \citep{Jimenez2025, TianReview2024}. To answer the research questions, three core policy scenarios were simulated by manipulating the DRL agent's available action space, using its core levers: Information, Education, and Communication (IEC), Incentives, and Enforcement (penalties).

To systematically evaluate the efficacy of the DRL agent, the simulation was structured to compare the \textit{Original}, \textit{Modified}, and \textit{New} policy frameworks:
\begin{itemize}
    \item \textbf{Original Policy (Baseline):} Corresponded to the \textit{Pure Penalty Regime}, representing the current ``No Segregation, No Collection'' approach with existing budget constraints.
    \item \textbf{Modified Policies (Static Scenarios):} Corresponded to the \textit{Pure Incentive} and fixed \textit{Hybrid Regimes}, representing manual attempts to alter implementation strategies without AI optimization.
    \item \textbf{New Policy (Optimized Output):} Corresponded to the final \textit{Adaptive Strategy} learned by the agent in the Hybrid setting, where the allocation of funds across the three levers dynamically shifted in response to household behavior.
\end{itemize}

\noindent The specific experimental regimes were defined as follows:

\begin{itemize}
    \item \textbf{Pure Incentive Regime:} Simulated a ``soft'' governance approach by limiting the DRL agent's action space to allocating funds only between IEC and Incentives (with no allocation for enforcement). This scenario tested the efficacy of relying solely on positive rewards, behavioral nudges, and social motivation to drive compliance \citep{LoanBalanay2023}. Research strongly supported that an incentive-based strategy could be effective in promoting pro-environmental behavior like waste separation \citep{Chen2023, MuZhang2021, Vorobeva2022}.
    
    \item \textbf{Pure Penalty Regime:} Simulated a ``hard'' governance approach by constraining the agent to allocating funds only between IEC and Enforcement. This scenario investigated the performance of a strategy heavily reliant on monitoring, fines, and institutional coercion \citep{Wang2023}. Given the challenges in enforcement highlighted in the context of the Ecological Solid Waste Management Act (R.A. 9003) in the Philippines, this scenario tested if a strictly punitive financial strategy could achieve superior compliance \citep{Dalugdog2021}.
    
    \item \textbf{Hybrid Regime:} Granted the agent access to its full action space (IEC, Incentives, and Enforcement). This scenario tasked the agent with finding the optimal, data-driven mix among all three policy levers. Literature suggested that an effective solution often required this sophisticated combination, using incentives to encourage behavior while employing penalties to deter non-compliance, demonstrating the superiority of a balanced, adaptive strategy over a single-instrument focus \citep{MuZhang2021, Zhao2022}. By comparing the compliance rates and cost-efficiency results across these three scenarios, the research could provide a powerful, data-driven recommendation for the LGU's long-term budget strategy \citep{TianReview2024}.
\end{itemize}

\subsection{Performance Metrics}
To evaluate the outcomes of the simulations and compare the success of each policy regime, four key performance metrics were defined, reflecting the multi-faceted nature of effective local governance.

\begin{enumerate}
    \item \textbf{Maximum Sustainable Compliance:} Defined as the highest population-weighted average compliance rate the LGU agent managed to achieve and maintain over the long term \citep{ApostolJamoralin2024, Torkayesh2021}. This established the agent's ability to solve the core environmental problem over a sustained period.
    
    \item \textbf{Cost-Effectiveness:} Calculated as the average percentage of compliance gained per \textpeso 100,000 spent. By integrating this financial ratio, the metric addressed the need for fiscal discipline and efficiency, ensuring the policy provided the most impact for its cost—a major practical concern for LGUs \citep{MedinaMijangos2020, WorldBank2022}.
    
    \item \textbf{Policy Equity:} Assessed by measuring the final variance in compliance rates between the 7 barangays \citep{Jimenez2025}. A lower variance indicated a more equitable and evenly distributed policy outcome, verifying that the DRL agent’s resource allocation strategy did not simply concentrate resources on easy-to-manage areas, but rather succeeded in raising compliance across all geographic units \citep{TianReview2024}.
    
    \item \textbf{Optimal Resource Allocation:} Analyzed for the comprehensive Hybrid regime as a key output. This metric detailed the final learned budget split, in both percentage (\%) and peso (\textpeso) terms, across the three policy levers (IEC, Enforcement, Incentives) and all 7 barangays. This provided the essential prescriptive, actionable intelligence, revealing the data-driven policy mix deemed optimal for the LGU's specific context \citep{MuZhang2021, Zhao2022}.
\end{enumerate}

\subsection{Global Sensitivity Analysis}
To validate the structural integrity of the model and ensure its findings were not artifacts of specific initialization values, a Global Sensitivity Analysis (GSA) was conducted using the method called \cite{Sobol2001}.

This rigorous validation technique decomposed the variance of the model's output (Global Compliance) into fractions attributable to specific input parameters. The analysis stress-tested the model by varying the five core behavioral parameters ($w_a, w_{sn}, w_{pbc}, c_{effort}, decay$) by $\pm 20\%$.

Two indices were calculated for each parameter:
\begin{itemize}
    \item \textbf{First-Order Index ($S_1$):} Measured the direct contribution of the parameter to the output variance.
    \item \textbf{Total-Order Index ($S_T$):} Measured the contribution including interactions with other parameters.
\end{itemize}